\section{段落}

\subsection{分段}

\begin{itemize}
  \item 最长用的分段方式是空行，一个空行或多个空行都可以
  \item 使用\verb|\par|命令
  \item \textbackslash \textbackslash 手动换行，不是分段。区别是分段会自动缩进，而手动换行不会自动缩进。
\end{itemize}

\subsection{缩进}

缩进由\verb|\parindent|控制。

中文文档类(如ctexart)自动设置段落缩进为两个汉字宽。

如果要在某一段开头临时禁用缩进，可在段前使用\verb|\noindent|命令。

如果要在本来没有缩进的地方使用缩进，可以用\verb|\indent|命令产生一个长为\verb|\parindent|的缩进。

\subsection{行距}

行距就是相邻两行文字之间的距离。在下面的图片中，两个下划线之间的长度就叫行距。

\begin{figure}[htpb]
    \centering
    \includegraphics[]{行距}
\end{figure}

行距的调节一般使用倍数，比如两倍行距。而“单倍行距”又根据字体、字号、软件的不同而改变（不同软件中有不同的定义，没有一个通用的值）。在LaTeX中，默认的单倍行距是字体字号的1.2倍。

如果要修改全局的默认行距，则在导言区加入“\verb|\linespread|”命令进行修改即可。比如想要修改成2倍行距，用如下代码修改即可：

\begin{lstlisting}
    \linespread{2}
\end{lstlisting}

如果只想对一小部分做更改，可以用大括号把这一段括起来，在调用\verb|\linespread|命令后，还要添加\verb|\selectfont|命令刷新行距信息。

\begin{lstlisting}
    {
    \linespread{2.0} \selectfont
    % 两倍行距的文字
    }
\end{lstlisting}

\verb|\linespread|命令不仅可以改变正文行距，同时也把目录、脚注、图表标题的行距也改了。使用 \verb|setspace| 宏包仅改变正文的行距，包括调整数学公式、浮动体的行距。

\begin{lstlisting}
    \usepackage{setspace}
    \singlespacing % 1倍字体高度
    \onehalfspacing % 1.5倍字体高度
    \doublespacing % 2倍字体高度
    \setstretch{1.25} % 任意行距
\end{lstlisting}

同时也可以在正文中进行局部环境设置

\begin{lstlisting}
    \begin{spacing}{2.0}
        % 这里是两倍行距
    \end{spacing}
\end{lstlisting}

\subsection{Latex与Word行距的对应关系}

同样是一倍行距，LaTeX看起来就是要比Word要紧密一些，原因就在于前面所说的LaTeX中，默认的行距是字体字号的1.2倍，而新版本的Word默认行距是字体字号的1.3倍。那怎么将LaTeX的行距设置得跟Word一样呢？很简单，就是“乘以1.3除以1.2”，比如在Word中要求的是1.5倍行距，则在LaTeX中设置行距为：
\[
    1.5 \times 1.3\div 1.2=1.625
\]

\subsection{段间距}

段与段之间的垂直距离，由变量\verb|\parskip|控制。

\verb|\parskip|的默认值是橡皮长度 0pt plus 1pt。

在中文排版中，常使用

\begin{lstlisting}
    \setlength\parskip{0pt}
\end{lstlisting}

把段间距定义为固定长度。


\section{调整水平间距}
\subsection{不可换行的水平间距}

\begin{itemize}
  \item \verb|\thinspace 或 \,| \qquad 0.1667 em
  \item \verb|\negthinspace| \qquad \quad \ -0.1667 em
  \item \verb|\enspace| \qquad \qquad \qquad 0.5 em
  \item \verb|\nobreakspace 或 ~| \quad 空格
\end{itemize}

\subsection{可换行的水平间距}

\begin{itemize}
  \item \verb|\quad| \qquad \ \ \ 1 em
  \item \verb|\qquad| \qquad \ 2 em
  \item \verb|\enskip| \qquad 0.5 em
  \item \verb|\|$\sqcup$  \qquad \qquad 空格
  \item \verb|\hspace{距离}| \qquad 产生指定的水平间距
  \item \verb|\hfill| \qquad 橡皮长度，\verb|\hspace{\fill}|的缩写，把内容均匀分布在一行
  \item \verb|\hspace{\stretch{2}}| \qquad 可设定填充的比例
  \item \verb|\hrulefill| \qquad 用横线填充
  \item \verb|\dotfill| \qquad 用圆点填充
  \item \verb|\phantom{}| \qquad 幻影盒子，产生与参数内容一样大小的盒子
\end{itemize}